% ---------------------------------------------------------------------------
% Related work for the switching-cost version, at two pages. The weighting differs from
% the inexpressibility version's: the attribute-grammar thread and the standardisation
% effort are load-bearing here, because this paper's result is that the translation needs
% attributes the query text does not carry, and the expressiveness literature is
% compressed to what §7 still needs.
% ---------------------------------------------------------------------------

\section{Background and Related Work}
\label{s2:sec:background}

Five traditions bear on the question of \S\ref{s2:sec:intro} and none of them answers it.

\paragraph{Query languages as formal languages, and where the attributes come from.}
Rich's treatment of query languages as formal languages is the licence for the approach,
and the regular/context-free/decidable classification is the book's frame rather than that
section's~\citep{rich2008}. An unambiguous grammar gives each query exactly one parse tree,
which is the precondition for a translation being well defined; an LALR(1) one makes the
parse linear and the parser a table, so a corpus can be processed in bulk rather than
sampled~\citep{aho2006}; and the translation itself has a classical form, a syntax-directed
transduction associating an output with each production and composing them over the parse
tree~\citep{lewis1968}.

One tradition in that list does more work here than the rest, and it is the one this
paper's result is about. Where a node's output depends on information held elsewhere,
Knuth's attribute grammars are the standard account~\citep{knuth1968}: an attribute is
computed and passed, and the grammar says where from. The translation this paper builds
needs two attributes --- the rollup interval and the metric's type --- that appear in no
production of either grammar and in no part of the query text, so there is nowhere in the
parse tree they could be inherited from. That is not a defect in the formalism but the
precise statement of \S\ref{s2:sec:unsaid}'s finding: the attributes exist, the translation
needs them, and the document that supplies them is the dashboard rather than the query.

\paragraph{The two languages.}
PromQL is documented and its reference parser is open source, so a grammar can be checked
against an implementation rather than inferred~\citep{promql}. It is a vector-valued
expression language in which selectors denote sets of series, range selectors denote
windows, functions reduce windows and aggregation operators reduce across series, with
counters stored as running totals and reset handling inside the range functions. Datadog's
metric query language is not documented as a grammar; what is published is a reference
manual with examples, enough to infer the syntax and not to verify it. The important
difference is not syntactic: Datadog stores a counter as per-interval deltas, and a query
is evaluated first by time and then by space~\citep{datadog,ddnested} --- a rollup reduces
each series alone and a space aggregator combines the results. Neither the rollup interval
nor the metric's type is part of the query text in either language, which is where this
paper ends up.

\paragraph{What a translation guarantees, and what the existing tools report.}
SQL dialect translation is the closest established practice: Apache Calcite parses SQL into
a relational algebra, optimises, and generates dialect-specific SQL on the way
out~\citep{begoli2018}, and SQLGlot does the same across more than twenty
dialects~\citep{sqlglot}. Both are engineering successes and neither offers a statement of
what a translation preserves --- their criterion is that the output parses and behaves as
users expect, and where a construct has no counterpart the usual outcome is a best-effort
rewrite that leaves no record of the approximation. That is a reasonable engineering choice
and a poor measurement instrument, which is why \S\ref{s2:sec:framework} requires a
translator to be total with every failure a named $\bot$.

Converters for this language pair exist, and we have found no published result among them
--- a negative about a literature, which is the weakest kind of claim this paper makes and
is stated as a search rather than as a fact. Chronosphere describes the closest thing to
this paper's method, a Datadog grammar written from scratch with queries parsed to an AST
and emitted through the open-source Prometheus grammar, and reports getting ``about
$90\%$'' of the way~\citep{marcopoli2024}; what the percentage is a percentage of, and what
a converted query is meant to preserve, are not stated. OpenObserve's dashboard migrator
marks each object migrated, needing review, or skipped~\citep{openobserve} --- a coarse
form of the typed failure \S\ref{s2:sec:framework} requires, reported without counts, and
the tool \S\ref{s2:sec:baseline} runs over this corpus. A third is closer to this paper
than either: a public repository with no accompanying publication, which translates Datadog
dashboards and monitors into PromQL and reports reaching $92.4\%$ of $14{,}172$ distinct
queries and monitors~\citep{dd2dash0} --- its own figures, read from a cached copy of its
README because the repository is a feature branch. It arrives independently at three of
this paper's devices: each rule carries one of four claims, each translated item carries a
ledger of the assumptions its claim rests on, and anything without a rule is refused by
the name of the construct rather than approximated. That convergence is the useful part of
finding it --- these are not this paper's inventions but what the problem forces on anyone
unwilling to round a translation up. What none of them publishes is a grammar for the
Datadog side, or a measurement of what is impossible as against unbuilt, or a count of the
context its output depends on and does not have.

\paragraph{The standard of proof, and the form of a negative.}
Compiler verification supplies the first: McCarthy and Painter proved a compiler for
arithmetic expressions correct by relating source and target
denotations~\citep{mccarthy1967}, and CompCert carries the idea through a realistic C
compiler with machine-checked proofs~\citep{leroy2009}. That is the shape of this paper's
soundness theorems though not their rigour, and it differs in one way that is not
incidental: a verified compiler may assume source and target run on the same input, and
here they do not. Database theory supplies the adjacent decision problem, equivalence of
conjunctive queries being decidable by the homomorphism
characterisation~\citep{chandra1977,abiteboul1995}; it does not transfer, since the
quantities here are numeric aggregates over time windows and exactness has to be stated up
to floating point, so the borrowing is the habit rather than the result. ``PromQL cannot say
this'' has a literature too~\citep{felleisen1991,mitchell1993}, from which this paper takes
the \emph{distinction} rather than the criterion: \S\ref{s2:sec:reverse}'s results are
blunter, containment facts about two classes of regular languages where the target has no
term with the denotation at all. They are the only claims here that hold regardless of how
much engineering is applied, which is exactly why the rest of the paper is careful not to
borrow their authority.

\paragraph{Observability, standardisation, and switching costs.}
The closest observability work compares implementations of \emph{one} language rather than
translating between two. Volz's compliance testing found that vendors advertising PromQL
support disagree with Prometheus on a generated conformance suite~\citep{volz2020} --- this
paper's problem in miniature, and the source of its practice of verifying a candidate
against a real server rather than a model. Standardisation approaches it from the other
end: OpenTelemetry has largely settled the \emph{data} question~\citep{otel}, leaving
queries as the remaining lock-in surface, and the CNCF's Observability TAG has an active
effort towards a common query language~\citep{sharma2025}. That effort needs exactly the
evidence this paper produces, and \S\ref{s2:sec:unsaid} has an uncomfortable thing to tell
it: agreeing on syntax and semantics for metric queries would leave most of the cost
measured here untouched, because what the corpus is missing is not in the query.

Switching costs themselves are well developed in industrial organisation. Farrell and
Klemperer survey lock-in and its competitive effects~\citep{farrell2007}, and Shapiro and
Varian give the lock-in cycle its standard informal treatment for information
goods~\citep{shapiro1999}. Two of their structural points land here: that what retains a
customer is not that leaving is impossible but that the terms of return differ from the
terms of entry, which \S\ref{s2:sec:reverse} is the test of; and that a switching cost need
not sit in the thing being switched but can sit in a complementary asset nobody priced,
which is what a resolution layer belonging to the dashboard tool turns out to be. The
concrete case that motivates the question is on the public record~\citep{orosz2023}.
